\chapter{Conclusão}
\label{cap:conclusao}

Ao longo destas semanas de trabalho, foi possível passar de uma representação
estática do problema para uma abordagem baseada em sensores simulados,
eventos, factos e regras capazes de acompanhar a evolução do ambiente ao longo
do tempo. A utilização do Drools permitiu pôr em prática esta abordagem através
de uma sessão com estado e um conjunto de regras capazes de produzir
diferentes propostas de decisão. A existência de mecanismos de prioridade,
resolução de conflitos e rastreabilidade permitiu acompanhar o processo de
inferência desde o início até à escolha da decisão a tomar.

O contacto com o perito teve também um papel importante na evolução do
trabalho, sobretudo na revisão de critérios relacionados com risco, leituras,
prioridades e comportamento perante situações de maior criticidade. Apesar de
ter ocorrido numa fase já avançada do desenvolvimento, permitiu identificar
aspetos que devem continuar a ser revistos e aprofundados.

O projeto encontra-se ainda em evolução, no qual existem componentes que
necessitam de ser desenvolvidos com maior pormenor e complexidade,
nomeadamente a implementação completa do \emph{Why Not}, a integração de
alguns critérios definidos no modelo conceptual, o tratamento mais completo da
informação sensorial e a conclusão da componente em Prolog. Estas limitações
não impedem, contudo, que o protótipo atual demonstre de forma clara a relação
entre aquisição de conhecimento, representação, inferência, tomada de decisão
e explicação.

De forma geral, o trabalho permitiu construir uma base consistente para
continuar o desenvolvimento do sistema, com o foco principal na forma como o
conhecimento pode ser estruturado, utilizado e analisado no processo de
decisão automatizado.
